\chapter{The Absence Of Truth In The Presence Of The Question}
\label{chap:absence-of-truth}

% Closed question: What remains when all the necessary names have been spoken?

The procession has reached its quiet final name. \lean{INFERRED} licenses
the public sentence the record carried itself through thirty-six gates to
earn. What it does not license is ownership of the truth the sentence
points at. The record permits the buyer to walk away from this car at
this price. It does not deliver the car's secret essence. The record
permits the controller to close the flight. It does not bind the flight's
next leg to obey.

Truth has not disappeared. It appeared in the procession at \lean{TRUTH}
and did real work there. The procession then kept moving through witness,
contact, locality, universality, calibration, halt, measurement,
compilation, and inference. Each later name was the next obligation truth
had to accept in order to remain public. The chapter that closes this
volume asks what the procession has left undone, why it had to be left
undone, and what kind of question replaces possession of truth when the
necessary names have all been spoken.

\section{Why \lean{TRUTH} Is Not The Ending}
\label{se:why-truth-is-not-the-ending}

The temptation is to end at \lean{TRUTH}. A reader is led through
admission, observation, comparison, and measurement, and then the final
word arrives: truth. The word seems to inherit everything that came
before it.

That transfer is dangerous. The discipline belongs to the procession,
not to a single banner word. Once a sentence is pulled away from the
record that carried it, it can travel as slogan. ``The car is worth
buying'' becomes ``that model is good,'' then ``those cars are safe
bets,'' then a neighbor's advice about a different car under different
conditions. The words have social force, but the record has been
stripped. The procession's discipline is precisely the discipline of
not letting the sentence be stripped.

\lean{TRUTH} is therefore not the ending. It is one obligation inside
the procession, important enough to be an anchor, but no more permitted
to travel alone than any of the other obligations. In Chapter
\ref{chap:argument} it named truth with a carriage history: a sentence
whose path can be checked, revised, or refused. The buyer's sentence is
honest only when it keeps that carriage. Under this inspection, at
this price, for this use case, and with these unresolved conditions
named, the buyer may infer worth. That is less grand and more
trustworthy.

An unhedged truth would be louder. It would also be less defensible.
If the car later fails in a way the inspection did not promise to
predict, the buyer who claimed unhedged truth would have to defend a
false kind of finality. The conditional inference avoids that failure.
It says exactly what the procession earned and no more. A speaker who
has crossed the thirty-six gates with the buyer cannot claim more than
the gates licensed; a speaker who has not crossed the gates is not the
buyer.

There is a deeper reason \lean{TRUTH} cannot be the ending, and this
volume is not the place to prove it. The selection rule that decides
which admissible continuation is the one actually carried forward
requires the kind of information measurement that no general-purpose
machine can compute. The compiler that built this project's
formalization arrives at that rule and cannot run it. The procession's
quiet final line is therefore not ``we have arrived at Truth'' but
``even the instrument that wrote this book cannot own the rule it just
named.'' The buyer's walking-away is the homely echo of that
cosmological refusal.

\section{The Question Remains}
\label{se:the-question-remains}

The final artifact is not a trophy. It is a disciplined question made
usable.

What can be distinguished? What can be admitted? What can be counted,
encoded, compared, represented, observed, held open, staged, executed,
scaled, witnessed, localized, halted, measured, compiled, and inferred?
Each name answers the question for the present record. None closes the
field forever. The next record will ask each of the names again.

For the car, the question remains in the most ordinary way. Tomorrow
the car will have new miles, new weather, new wear, new risks. A month
from now the same buyer may need another inspection, or the same car
may have become a different car under a different owner. The earlier
inference was not false because new facts arrived. It was current under
the record that licensed it. The new record may license a different
sentence, or the same sentence, or no sentence at all; the procession
will decide when it is run again.

For the flight, the same shape applies. The flight closed under the
record. The next flight is a different flight. The aircraft that
landed clean may take off again with the same registration and a
different fuel load, a different crew, a different weather. The
procession will run again. The clearance the controller issued tomorrow
is a new clearance, not a continuation of yesterday's.

A gauge does not own the thing it reads. A speedometer does not own
speed. It produces a reading under calibration, and the reading is
current until the next reading. The next moment may need another
reading. The rhetorical gauge is the same. It gives a speaker
permission to speak under what the record has admitted, and the
permission is real, and the permission is limited, and the limitation
is part of the permission.

This is why the book has followed names rather than trying to define
Truth from above. The names keep the question alive. They let a public
sentence be spoken without pretending that the world has stopped
changing underneath it. A book that ended at Truth would have to claim
that the naming was complete. A book that ends at the question can be
honest about the next reading the next instrument will ask for.

The procession is a method, not a possession. The method can be run by
anyone who has read the names, against any record the procession
admits. The method does not promise that the same record run twice
will yield the same sentence; it promises that the differences will be
inspectable.

\section{Inference Without Ownership}
\label{se:inference-without-ownership}

\lean{INFERRED} is weaker than ownership and stronger than opinion.

It is weaker than ownership because the speaker does not possess the
future. The buyer does not own the transmission's next failure, the
road salt's slow work, or next year's repair market. The speaker has a
conclusion about the present record and the present question. The
conclusion may be acted on. The action is itself a record, and that
record will enter the next procession.

It is stronger than opinion because the sentence is licensed by a
public procession. A reader can challenge the inference by challenging
particular links: the comparable set, the repair estimate, the title
gate, the calibration of the threshold, the missing receipt. The
disagreement has handles. That is what opinion lacks. An opinion can
be refused only as a whole; an inference can be revised at the
specific gate where the prior record failed.

The space between ownership and opinion is where honest public
argument lives. Ownership offers false finality. Opinion offers private
license. Inference offers neither comfort. It requires a speaker to
walk the procession and then admit what the procession does not cover.

That is not a weakness in the project. It is the point of the
rhetorical gauge. The names do not give the speaker the world. They
give the speaker permission to speak under the record. The buyer may
buy, renegotiate, or walk away without claiming the car's essence. The
controller may close the flight, file an incident report, or release
the aircraft to maintenance without claiming that the flight has
predicted the next flight. Both moves are inference without ownership:
a public action licensed by a record and bounded by a question.

The same shape applies to the project's largest claim. The rule that
selects which admissible continuation is the one carried forward needs
a measurement of informational content that no Turing procedure can
perform exactly. Lean itself, the compiler that proved every other
obligation in this book, arrives at this rule and cannot run it. The
manga that accompanies this project ends with a blinking cursor over
the word \lean{sorry}. This is not failure either. It is the same
inference-without-ownership played at the scale of the universe: the
system may name the rule it cannot own. Naming the rule is the gauge's
last honest move.

\begin{closingmotion}

The rhetorical gauge ends where the reference trace can begin again.

The names made the argument public. The technical volumes explain why
each name carries the obligation the name claims. The reference trace
asks Lean what it accepts and what it refuses. The historical-experiment
volume shows where each name first earned itself in public. The manga
shows what it felt like to build.

Together they form one project. This volume carries the public sentence
the rest of the project licenses: the record does not give us Truth to
own. It gives us inference under names. Tomorrow's record will ask the
procession again, and the procession will be ready to run.

\end{closingmotion}
